\section{Data and Task Definitions}
\label{sec:data}

\subsection{PTB-XL}
We used PTB-XL v1.0.3~\cite{wagner2020ptbxl}, a public collection of 21,799 ten-second, 12-lead resting ECGs from 18,869 patients, annotated by cardiologists with SCP-ECG statements. Following the reference implementation of Strodthoff et al.~\cite{strodthoff2021benchmark}, every diagnostic statement was mapped to one of five diagnostic superclasses---normal ECG (NORM), myocardial infarction (MI), ST/T change (STTC), conduction disturbance (CD) and hypertrophy (HYP)---regardless of its likelihood value; the 21,388 records with at least one superclass were kept. We used the 100\,Hz version (1,000 samples per lead), as in the PTB-XL benchmark; it keeps the waveform content needed for superclass diagnosis while keeping the full experiment grid of more than 100 training runs within the budget of one laptop GPU.

We used the official stratified folds without modification: folds 1--8 for training (17,084 records), fold 9 for validation (2,146) and fold 10 for testing (2,158). The folds keep all records of a patient in one fold, so there is no patient-level leakage. Label prevalence is almost identical across splits (\cref{tab:data}). Signals were z-scored per lead with a mean and standard deviation estimated on the training fold only, and the same transform was applied to validation, test and external data.

\subsection{Tasks}
\textbf{T1 -- multi-label diagnosis.} Five sigmoid outputs, one per superclass; a record may carry several labels.

\noindent\textbf{T2 -- normal-versus-abnormal triage.} One sigmoid output; a record is \emph{abnormal} if it carries any of MI, STTC, CD or HYP. T2 is the simpler question a wearable is most likely to be asked: does this person need a 12-lead ECG? T2 is not the complement of NORM: 51 test records are labelled NORM together with an abnormal superclass, and they count as abnormal in T2. For T1 models we also report a \emph{derived} triage score, $1-P(\text{NORM})$, evaluated against the same T2 label.

\subsection{Lead configurations}
Each model was trained and tested on one of six lead configurations chosen to span clinical and wearable devices (\cref{tab:leads}).

\begin{table}[H]
\centering
\small
\caption{Lead configurations.}
\label{tab:leads}
\begin{tabular}{lll}
\toprule
Configuration & Leads & Rationale \\
\midrule
12lead & all 12 & clinical reference \\
limb6 & I, II, III, aVR, aVL, aVF & limb-only devices, six-lead patches \\
chest6 & V1--V6 & precordial information without limb leads \\
2lead (I + II) & I, II & two-electrode devices \\
leadII & II & classic rhythm strip \\
leadI & I & smartwatch-like single lead \\
\bottomrule
\end{tabular}
\end{table}

\subsection{External cohort}
For external validation (E7) we used the Chapman-Shaoxing/Ningbo database~\cite{zheng2022physionet,zheng2020chapman} (45,152 ten-second 12-lead ECGs at 500\,Hz from two Chinese hospitals), whose diagnoses are SNOMED-CT codes. We mapped codes to the PTB-XL superclasses with a fixed table (``v2 mapping'') that follows the official equivalence list of the PhysioNet/CinC Challenge 2021~\cite{reyna2021willtwodo}; for example, complete right bundle branch block is treated as right bundle branch block. A record was \emph{normal} only if all its codes were sinus-rhythm codes (sinus rhythm, sinus bradycardia, sinus tachycardia, sinus arrhythmia), and \emph{abnormal} if at least one code mapped to MI, STTC, CD or HYP. Records with only unmapped findings (for example atrial fibrillation alone, axis deviation or ectopic beats) were excluded, because PTB-XL's diagnostic superclasses do not cover them. Signals were converted to millivolts, resampled to 100\,Hz with an anti-aliasing polyphase filter, cropped or zero-padded to 10\,s, and records with a flat or non-finite lead were excluded. Of 45,152 records, 34,815 were included (16,092 abnormal, 18,723 normal); 8,971 were excluded for unmapped diagnoses and 1,366 for signal problems (\cref{tab:data}). Two alternative mappings were analysed as sensitivity analyses (\cref{sec:res_triage}).

\begin{table}[H]
\centering
\small
\caption{Cohorts. PTB-XL label prevalence is multi-label (rows do not sum to 100\%). Chapman fractions use all scanned records.}
\label{tab:data}
\begin{tabular}{lrrrrrr}
\toprule
PTB-XL split & Records & NORM & MI & STTC & CD & HYP \\
\midrule
Train (folds 1--8) & 17,084 & 44.5\% & 25.6\% & 24.5\% & 22.9\% & 12.4\% \\
Validation (fold 9) & 2,146 & 44.5\% & 25.2\% & 24.6\% & 23.1\% & 12.5\% \\
Test (fold 10) & 2,158 & 44.6\% & 25.5\% & 24.1\% & 23.0\% & 12.1\% \\
\midrule
\multicolumn{7}{l}{\textbf{Chapman-Shaoxing/Ningbo (external, v2 mapping)}} \\
\multicolumn{5}{l}{Scanned} & \multicolumn{2}{r}{45,152 (100\%)} \\
\multicolumn{5}{l}{Excluded: no mappable diagnosis} & \multicolumn{2}{r}{8,971 (19.9\%)} \\
\multicolumn{5}{l}{Excluded: flat lead / non-finite signal / unreadable header} & \multicolumn{2}{r}{1,290 / 75 / 1} \\
\multicolumn{5}{l}{Included (abnormal / normal)} & \multicolumn{2}{r}{34,815 (16,092 / 18,723)} \\
\bottomrule
\end{tabular}
\end{table}
